% Part: lambda-calculus
% Chapter: lambda-definability
% Section: truth-values

\documentclass[../../../include/open-logic-section]{subfiles}

\begin{document}

\olfileid[fr]{lam}{rep}{tvr}
\olsection{Valeurs de vérité et relations}

Dans le lambda-calcul pur, on peut coder
les valeurs de vérité par
\begin{align*}
  \fn{true} & \ident \lambd[x][\lambd[y][x]]\\
  \fn{false} & \ident \lambd[x][\lambd[y][y]].
\end{align*}

Ces valeurs sont des \emph{sélecteurs} :
elles reçoivent deux arguments et
renvoient l'un d'eux. Le terme
$\fn{true}$ sélectionne le premier,
$\fn{false}$ le second. Ainsi,
$\fn{true}\,M N\bred M$ et
$\fn{false}\,M N\bred N$.

\begin{defn}
Une relation $R\subseteq\Nat^k$ est
\emph{$\lambda$-définissable} s'il existe
un lambda-terme~$T_R$ tel que
\begin{align*}
  T_R\,\num{n_1}\dots\num{n_k}
    & \bred \fn{true}
  \intertext{lorsque $R(n_1,\dots,n_k)$
    est vraie, et}
  T_R\,\num{n_1}\dots\num{n_k}
    & \bred \fn{false}
\end{align*}
sinon.
\end{defn}

Par exemple, la relation $\{0\}$,
vraie exactement pour~$0$, est
$\lambda$-définie par le terme
\[
  \fn{IsZero} \ident
  \lambd[n][n (\lambd[x][\fn{false}])\,\fn{true}].
\]
En effet, un numéral de Church est un
itérateur : nous partons de $\fn{true}$
et chaque itération renvoie $\fn{false}$,
quel que soit le résultat précédent.
Si $n=0$, aucune itération n'a lieu et
le résultat est $\fn{true}$; si $n>0$,
le résultat est $\fn{false}$.

Avec ce codage, nous pouvons définir
des fonctions de vérité. Voici la
négation et la conjonction :
\begin{align*}
  \fn{Not} & \ident
    \lambd[x][x\,\fn{false}\,\fn{true}]\\
  \fn{And} & \ident
    \lambd[x][\lambd[y][xy\,\fn{false}]].
\end{align*}
La fonction $\fn{Not}$ renvoie
$\fn{true}$ sur l'entrée $\fn{false}$,
et $\fn{false}$ sur l'entrée $\fn{true}$.
La fonction $\fn{And}$ reçoit deux
valeurs de vérité $x$ et~$y$, et doit
renvoyer $\fn{true}$ exactement lorsque
les deux sont $\fn{true}$. Comme $x$
est un sélecteur, l'application
$xy\,\fn{false}$ renvoie~$y$ si
$x=\fn{true}$, et $\fn{false}$ si
$x=\fn{false}$. C'est bien la
conjonction voulue.

Nous supposons ici que les arguments de
$\fn{And}$ sont des valeurs de vérité.
Si l'on lui passe d'autres termes,
le résultat, c'est-à-dire sa forme
normale lorsqu'elle existe, peut ne
pas être une valeur de vérité.

\begin{prob}
Définir $\fn{Or}$ et $\fn{Xor}$,
qui représentent respectivement la
disjonction inclusive et la disjonction
exclusive avec ce codage des valeurs
de vérité en lambda-termes.
\end{prob}

Pour représenter une relation à $k$ arguments, on applique
son terme représentant aux $k$ numéraux de l'entrée.
Par exemple, $\fn{IsZero}$ reçoit un seul numéral.
Le terme $\fn{IsZero}\,\num{0}$ se réduit à $\fn{true}$.
Pour $n>0$, le terme $\fn{IsZero}\,\num{n}$ se réduit à
$\fn{false}$.

\end{document}
